\documentclass[11pt]{article}
\usepackage[margin=1in]{geometry}
\usepackage{amsmath,booktabs}
\usepackage{xurl}
\usepackage[hidelinks]{hyperref}
\setlength{\emergencystretch}{2em}
\title{Multi-topic Modelling of LLM Chat Sessions for Guided Agent Search}
\author{Authorship to be finalized}
\date{Living research draft --- 1 October 2026}
\begin{document}
\maketitle
\begin{abstract}
We investigate how to organize histories of interactions with coding agents into
overlapping, interpretable topics that can guide subsequent retrieval. Long
sessions contain multiple subjects, while short replies may depend on context.
The intended system must update incrementally, preserve source provenance, and
support topic exploration and filtering without replacing ordinary search.
This draft records the research protocol and initial exploratory experiments:
twelve geometry controls on 664 independent threads and a small user-turn
boundary pilot. These experiments expose tradeoffs in order sensitivity,
abstention, broad clusters and oversegmentation. They do not yet establish a
selected method, reliable topic quality, or improved agent search performance.
\end{abstract}

\section{Problem and research questions}
A coding-agent session is not necessarily one document about one subject. A
session may discuss implementation, diagnosis, tooling and unrelated planning.
Conversely, a short approval can continue a substantive earlier discussion while
carrying little standalone lexical evidence. Our target is a source-grounded
collection of topic tags per thread, with meaningful ancestor inheritance and
tools for browsing topics, finding candidate threads, and searching messages
within those threads. A topic's discovery structure and its semantic taxonomy
are distinct objects; an internal clustering node is not automatically a useful
parent label.

We ask: (RQ1) which conversation units preserve primary and secondary subjects;
(RQ2) which discovery and representation methods produce specific, accurate
overlapping tags; (RQ3) how quality changes with insertion order, incremental
updates and provider; and (RQ4) whether topic guidance improves retrieval success
and cost for agents compared with the existing hybrid search workflow. These
are research questions, not claims already demonstrated. Novelty remains to be
established through a fuller literature review.

\section{Related work}
DevGPT provides developer--ChatGPT conversations linked to software artifacts
\cite{devgpt}. Its publicly shared conversations are useful external material,
but their selection and tool-trace coverage differ from local coding histories.
Contextual topic classification uses conversation history and dialogue-act
features \cite{khatri}; this motivates context experiments rather than wholesale
copying of assistant answers into every short user turn.

Utterance-pair coherence models \cite{xing} and topic-aware representations
\cite{gao} motivate evaluating dialogue episodes before clustering entire
sessions. These learned approaches are not equivalent to a threshold on raw
adjacent embedding similarity. BERTopic separates embedding-based discovery
from class-based TF-IDF representation \cite{bertopic}; we accordingly evaluate
label quality independently of discovery. Automatic coherence can disagree with
human judgments \cite{hoyle}, so it is not our sole selection criterion.

GRINCH uses rotation and grafting to revise an incremental hierarchy
\cite{grinch}. Its theoretical order-independent recovery requires assumptions
about linkage consistency; it does not establish accuracy on our embeddings.
MemTree studies dynamic hierarchical memory \cite{memtree}. Its memory-retrieval
evaluation is adjacent evidence, not validation of granular overlapping topic
tags. Our initial BIRCH implementation uses the scikit-learn reference API
\cite{sklearn}; the production choice remains open.

\section{Experimental protocol}
\subsection{Research sequence and evidence separation}
The agreed workflow is literature compilation, approach shortlist, Python
experiments with recorded results, method selection, then Rust implementation.
The Python environment is isolated from the Rust application. Private source
texts, native IDs and detailed assignments remain outside the Git repository;
this manuscript contains aggregate results. An existing Rust prototype is not
evidence that a topic feature is production ready.

Development annotations and model selection must be separated from reserved
evaluation. We intend to annotate primary and secondary subjects, irrelevant
framing, context-dependent turns, episode boundaries and semantic parents,
independently of candidate model outputs. Annotation uncertainty, incomplete
source context and disagreements must be retained. Independent human review and
inter-annotator agreement have not yet been completed. The boundary pilot below
uses a single coding-assistant reviewer and is not a human ground-truth study.

\subsection{Frozen data and reproducibility}
The initial local snapshot contains 789 native threads from Codex and Claude
Code, with 6,714 compatible stored user-message vectors. The controls remove 242
duplicate searchable user texts within threads and select one representative
for each canonical source group, removing 121 duplicate threads. The resulting
cohort contains 664 independent threads. A frozen partition report supplies
generation and canonical identities, not cleaned topic vectors.

The reserved partition consists of groups whose first identity byte modulo five
equals zero: 126 threads; held-out discovery uses the other 538. This partition
is fixed before fitting. Full discovery is an exploratory control, and reserved
coverage after full discovery is not a generalization measurement. Future
independent test data may be necessary because the corpus has also been used in
earlier exploratory development.

Stored vectors are 384-dimensional signed-byte representations. Each vector is
converted to floating point and normalized; each thread is represented by the
normalized mean of its unique user vectors. No model inference occurs in these
controls. The loader rejects incompatible dimensions/generations, missing groups,
invalid vectors, inconsistent persisted norms and a live SQLite WAL. It opens the snapshot read-only and hashes
the source fields, vector bytes, generation and exact partition report for
comparison integrity. The generation digest matches the BLAKE3 of the Rust CPU
specification: FastEmbed 6.0.1, AllMiniLML6V2Q, batch size 8, eight workers with
two threads each, per-vector symmetric signed-byte quantization and schema 3.
Rust scales each float vector's maximum absolute component to 127, rounds and
clamps components to $[-127,127]$, and stores their byte bit patterns and norm.
Python interprets these as signed bytes. Reconstructing the float32 norms matches
all 6,714 stored user norms exactly. Installed associated ONNX/tokenizer artifacts
were also hashed privately, but the generation specification does not hash model
bytes or prove historical asset identity. Complete model provenance and external
vector regeneration remain reproducibility tasks.

Runs use Python 3.11.15, NumPy 2.4.6 and scikit-learn 1.9.1 on the Linux Xenia
host (KVM guest, 16 visible vCPUs, AMD EPYC 9645 reported by \texttt{lscpu},
approximately 62.8 GiB guest RAM), with locked dependencies and one numeric CPU
thread. These guest specifications do not imply dedicated physical CPU access.
Each run starts a
fresh process. Fit and label times exclude loading and interpreter startup;
peak RSS includes them. Load isolation and broader hardware characterization remain
to be added. No repeated-run latency distribution or confidence interval is
claimed from a single forward run.

\subsection{Initial geometry and label controls}
We run three methods in full/held-out discovery and forward/reverse native
thread order, giving twelve runs. BIRCH uses threshold 0.4, branching factor 64
and no final global clustering. HDBSCAN uses Euclidean distance, minimum cluster
size 5, minimum samples 5 and excess-of-mass selection. Average linkage uses
cosine distance with cut threshold 0.3. Parameters are initial controls, not
tuned to equivalent semantic quality. HDBSCAN and average linkage are offline
comparators, not direct solutions to incremental maintenance.

For the label control, we concatenate training texts by fitted cluster, count
English-stopword-filtered unigrams and bigrams (maximum vocabulary 20,000), and
rank terms using
\begin{equation}
 s_{t,c}=\frac{f_{t,c}}{\max(L_c,1)}
          \log\left(1+\frac{\overline{L}}{\max(f_t,1)}\right),
\end{equation}
where $L_c$ is counted cluster length, $\overline{L}$ its mean across clusters,
and $f_t$ total counted training frequency. We retain the top eight terms.
Stopwords affect labels only, not stored embedding inputs. This is a c-TF-IDF
control, not a full BERTopic pipeline or a validated naming method.

We also record a common nearest normalized-centroid assignment at cosine
similarity at least 0.7. This is an uncalibrated diagnostic, not native HDBSCAN
prediction, overlapping assignment, or semantic correctness. Fit noise remains
distinct from diagnostic centroid assignment.

\section{Results and analysis}
\subsection{Geometry controls}
Table~\ref{tab:geometry} records forward-run counts/times and forward/reverse
coassignment Jaccard. The latter compares sets of within-cluster training pairs,
excluding noise. Empty pair unions must not be scored as perfect stability.
Cluster size and noise qualify interpretation: broad merging or abstention can
produce favorable stability without useful topics.

\begin{table}[ht]
\centering\small
\begin{tabular}{llrrrrrr}
\toprule
Method & Scope & Clusters & Noise & Largest & Fit ms & Label ms & Jaccard\\
\midrule
BIRCH & Full &166&0&104&30.8&1326.2&0.701\\
BIRCH & Held-out &141&0&101&34.5&1202.3&0.933\\
HDBSCAN & Full &40&225&36&192.7&896.3&0.983\\
HDBSCAN & Held-out &31&205&34&115.6&868.6&0.953\\
Average-link & Full &112&0&191&62.9&1254.7&1.000\\
Average-link & Held-out &100&0&155&48.7&1232.3&1.000\\
\bottomrule
\end{tabular}
\caption{Exploratory raw thread-mean controls, not quality-ranked methods.}
\label{tab:geometry}
\end{table}

Forward peak process RSS ranges from 228.7 to 237.3 MiB. Held-out centroid
coverage is 119/126 for BIRCH, 103/126 for HDBSCAN and 118/126 for average linkage.
Coverage is not assignment accuracy. BIRCH's fast fit comes with order sensitivity
in the full control. HDBSCAN's smaller clusters and substantial noise require
checking rare-topic recall; average linkage's 191-member largest cluster needs
semantic review. Labeling costs exceed fitting in these runs. This observation
does not imply optimal label implementation or end-to-end indexing latency.
GRINCH was not part of these twelve controls, and no winner is selected.

\subsection{Earlier embedding-cleanup pilot}
Before establishing the Python lab, a paired local-embedding pilot compared raw
and conversationally cleaned inputs on two synthetic corpora (25 documents and
21 positive queries each) and a source-reviewed short-request corpus (17 targets
plus 100 provider-balanced distractors). Both singleton and document-batch
embedding policies were checked. The real requests ranked first in all 17 cases
under both input policies, leaving no measured improvement or headroom. Synthetic
cases showed no aggregate gain and some near-tied ranking regressions. These
measurements used dense quantized retrieval rather than the complete hybrid
search pipeline, lacked independent exhaustive negative judgments, and cannot
establish that cleaning improves either topics or end-to-end retrieval.
Production embedding inputs remain unchanged. The detailed historical protocol
and results are recorded in
\path{benchmarks/2026-10-01-embedding-cleanup.md}; they are exploratory evidence,
not part of the twelve Python geometry controls.

\subsection{Exploratory user-turn boundary pilot}
We select the first two eligible development threads per provider in native-ID
order, requiring 6--20 nonempty searchable user turns and at most 1,500 characters
per turn. This purposive bounded selection yields four threads, 33 turns and
29 boundaries. It excludes reserved groups and is not representative random
sampling. Empty search projections are omitted; nonempty unembedded turns are
retained, and cosine predictors abstain on boundaries adjacent to missing
vectors rather than bridging them silently.

A single coding assistant reviews the user-only text sequence and records same
subject, changed subject or uncertain judgments with rationales before running
the current predictors. SHA256 pins the exact IDs, ordinals and searchable texts.
Six boundaries are uncertain, leaving 23 judged boundaries with six changes.
The reviewer has prior exposure to parts of the corpus; this is neither blinded
human annotation nor an independent test set. Assistant responses and native
branch structure are not reviewed here, so these judgments apply only to linear
user-text presentation.

Controls split every three user turns, split adjacent vectors below cosine 0.55,
or split against the current segment's normalized mean below 0.55. The mean resets
at a split or missing-vector boundary. The fixed-window predictor detects 2/6
changes with five false positives (precision 0.286, recall 0.333). Both cosine
controls detect 6/6 with fourteen false positives (precision 0.300, recall 1.000).
No predictor abstains on a judged boundary in these selected cases. Uncertain
judgments are excluded, not counted as correct predictions.

At this fixed threshold, the cosine controls oversegment the reviewed user-only
continuations. This motivates context-aware episode representations and boundary
evaluation; it does not prove a learned segmentation method is better or justify
a threshold selected on these examples. Label granularity itself needs clearer
guidelines and independent review before scaling the study.

\subsection{Bounded contextual boundary comparison}
This development comparison holds the existing annotations and cosine
threshold fixed. It compares adjacent vectors against two disjoint blocks of up
to two user turns on either side of each boundary, normalizing each block mean.
The right block uses at most one future user turn; decisions therefore require
bounded lookahead or provisional online publication. Blocks share no user turn:
shared-input smoothing could otherwise inflate similarity by construction.
If any vector in either block is missing, or a block mean cancels to numerical
zero, the method abstains. At thread edges,
available shorter blocks are used. This pooling control does not re-embed context,
resolve pronouns, inspect native branches or incorporate assistant text.
Its hypothesis is reduced spurious splitting at the risk of merging genuine
subject changes. On the same 23 judged boundaries, it detects all six changes
with sixteen false positives (precision 0.273, recall 1.000), compared with
fourteen false positives for adjacent/segment-centroid controls. No judged
boundary requires abstention in this pilot. Pooling therefore does not improve
boundary precision at this fixed threshold. We did not tune the threshold or
change annotations in response. Baseline results and input/annotation hashes
were verified unchanged. These fixed development choices are not a preregistered
independent evaluation; this negative result motivates testing genuinely
context-aware representations rather than assuming vector averaging resolves
referential replies.
The next annotation protocol is documented in
\path{research/topics/annotation-guidelines.md}; it specifies episode granularity,
source/context scope, secondary subjects, ancestors and disagreement handling.
It does not retroactively alter these exploratory judgments.

\subsection{Re-embedded text context control}
We next compare actual text concatenation with vector pooling, using the
explicitly installed MiniLM ONNX graph and tokenizer. The Python lab adds locked
ONNX Runtime 1.30.0 and Tokenizers 0.23.2 dependencies; it does not download
models. Inference/session and tokenizer APIs follow their official documentation
\cite{onnxruntime,tokenizers}. Singleton CPU inference uses one intra-/inter-op thread, token-mean
pooling with all singleton tokens attended, and per-vector symmetric signed-byte
rounding (ties away from zero). We reject inputs exceeding 512 tokens. A bounded
cache avoids repeated identical texts within the run. Disjoint blocks contain
up to two user turns joined by two newlines and need one-user lookahead. This
does not resolve references through assistant answers or native branches.

We retain the same source/annotation fingerprints and cutoff 0.55. A re-embedded
single-turn control detects six changes with fifteen false positives (precision
0.286). Re-embedded text blocks detect six with seventeen false positives
(precision 0.261): every judged continuation is split. Both have recall 1.000
without abstentions on judged boundaries. These results do not demonstrate an
improvement; concatenation alone is not a validated dialogue representation.

Across the 33 single turns, stored/re-embedded cosine ranges from 0.9654 to 0.9946
(median 0.9888). The extra false positive in the singleton control shows that
runtime/batch/input-policy differences can affect a fixed cutoff. We do not claim
byte-identical parity with historical Rust inference. The installed asset hashes
are \path{afdb6f1a0e45b715d0bb9b11772f032c399babd23bfc31fed1c170afc848bdb1}
for the ONNX graph and
\path{da0e79933b9ed51798a3ae27893d3c5fa4a201126cef75586296df9b4d2c62a0}
for the tokenizer. Generation identity still does not prove historical model
bytes or a fully reproduced embedding procedure.

The exploratory run made 62 unique singleton inference calls and 87 cache hits;
maximum input length was 55 tokens, so no truncation occurred. Inference took
120.5 ms, tokenization 7.7 ms, and measured elapsed loading/experiment time
662.1 ms; process peak RSS was 751.2 MiB. Elapsed time begins after Python imports,
while RSS includes them. These single-run values are not latency distributions,
an optimized online implementation, or measurements of the Rust local worker.
The experimental report is \path{text-context-results.json} in the private
boundary-pilot artifact directory. Subsequent reports also pin source/lock hashes.
An independent fresh-process repeat reproduced predictions, drift values,
fingerprints and inference/cache counts exactly, excluding time/RSS measurements;
every compared context had disjoint component message IDs.
The re-embedded controls strengthen the negative result at this cutoff but do
not establish whether learned dialogue coherence, calibrated boundaries or
source-grounded assistant context will perform better.

\subsection{Expanded development source review}
We add eight development conversations, excluding the first pilot's canonical
source groups and all reserved groups. Selection takes two groups per provider
and length stratum: 3--10 or 11--40 nonempty searchable user records. Each record
is at most 4,096 characters and each thread at most 40,000 UTF-8 bytes. A fixed
seed, \texttt{mash-topic-review-v2-20261001}, ranks each stratum by SHA256 of seed
bytes followed by canonical group bytes, independently of model outputs.
The eligible pool contains 21 groups: Claude Code medium/short 4/3, Codex
medium/short 6/8. These restrictive bounds and a single user's archive introduce
substantial selection bias; this is not a representative corpus study.

The selected sequence contains 84 visible user-role records and 76 boundaries.
Before predictions for this sample, one coding assistant freezes admission,
subject and boundary judgments from searchable user texts under the annotation
guidelines. The review identifies 39 substantive-content records, 29 referential
replies, 13 protocol/generated records, two attachment-only records and one
redacted credential-like value. Role alone therefore does not identify substantive
user intent. An opaque-value display guard preserves source hashes but is not a
general secret scanner; private review artifacts must not be printed or released
without further review. These are manual source judgments, not a production
admission classifier or complete native-envelope verification.

The subject sidecar records nine primary, nine secondary and three uncertain
subject judgments with supporting native ordinals. Six of eight conversations
have multiple source-supported subjects. This is evidence to retain overlapping
tags in the evaluation design, not a measured recall score or validated taxonomy.
The user-only boundary review leaves 36 boundaries uncertain, including framing,
unresolved referents, unavailable attachment content and ambiguous granularity;
40 boundaries remain known, with twelve changes. We retain those uncertainties
rather than treating them as correct continuation predictions.

\begin{table}[ht]
\centering
\begin{tabular}{lrrrr}
\toprule
Control & True changes & False splits & Precision & Recall\\
\midrule
Window of three &3&9&0.250&0.250\\
Adjacent cosine 0.55 &12&26&0.316&1.000\\
Segment-centroid cosine 0.55 &12&26&0.316&1.000\\
Disjoint block means 0.55 &12&25&0.324&1.000\\
\bottomrule
\end{tabular}
\caption{Unchanged-parameter controls on the expanded development review;
36 uncertain boundaries excluded. No judged boundary requires abstention.}
\label{tab:expanded}
\end{table}

Table~\ref{tab:expanded} again exposes oversegmentation at the uncalibrated
cutoff, rather than establishing useful episode quality. Of known boundaries,
Claude Code contributes twelve with three changes; Codex contributes 28 with
nine. Adjacent-control precision is 0.273 and 0.333 respectively. The small
one-false-split difference for block means does not establish a reliable gain.
No parameters or labels were adjusted after these predictions.
All groups are distinct development groups; input/source pins, complete admission
IDs and subject-support ordinals were checked. The review remains single-assistant,
user-only and familiar with the broader archive, not independently annotated
human ground truth. Full conversational and native-branch context is required
before resolving uncertainties or selecting a production topic method.

Private artifacts are recorded under \path{topic-python-review-v2-20261001}:
\path{packets.json}, \path{annotations.toml}, \path{source-review.json} and
\path{boundary-baselines.json}. The sampling implementation is
\path{research/topics/review_packets.py}. Earlier pilot annotations remain
unchanged, and these development results do not consume the reserved partition
as a tuning set or establish external generalization.

\section{Planned evaluation and unresolved requirements}
Next experiments will compare bounded dialogue units and source-grounded context
selection, holding discovery fixed; then label-only contrastive training-background
weights and representative-phrase scoring, holding memberships fixed. Online
discovery alternatives come after unit/label evaluation. Changes to geometry
must be frozen/versioned so previously inserted statistics remain compatible.
Edits, deletions, late replies and provisional episode boundaries require explicit
lifecycle semantics and measured incremental costs.

The topic evaluation will report primary/secondary assignment precision and
recall, rare-topic coverage, label specificity/intrusion, semantic parent
correctness and inherited ancestor correctness, provider breakdowns, order
stability, wall time, RSS and inference cost. Future guided-search tasks will
compare direct hybrid search, thread-first navigation and topic-guided navigation
under equal query/tool budgets, reporting relevant-source success, ranking and
time/cost with uncertainty estimates. These tasks have not yet been conducted.

Private histories limit release and external reproducibility. Before a public
evaluation, establish reviewed release permissions or an appropriate external
corpus, record complete model/hardware provenance, obtain independent annotations,
and preserve development/test separation. Current geometry and boundary results
are exploratory diagnostics. Robust overlapping tags, validated semantic trees,
MCP topic tools and improved downstream agent search remain unproven.

\bibliographystyle{plain}
\bibliography{references}
\appendix
\section{Experiment ledger}
The geometry controls are implemented in \path{research/topics/run.py}; the
boundary pilot in \path{research/topics/episodes.py}. Dependencies are pinned in
\path{research/topics/uv.lock}. Private artifacts are recorded under
\path{topic-python-baselines-20261001} and
\path{topic-python-boundaries-20261001}. The source snapshot is
\path{search-performance-d62a597/index.sqlite3}; the frozen partition report is
\path{topic-option-context-clusters-20261001/control-birch-full/UserRepresentatives-0.4-false.json}.
These artifacts are not publicly released data.

\begin{description}
\item[Embedding generation] \path{ca9c7766f955402eb34db3c324da24a512e1118de44e3d9aa9a0faeeac0da5f4}
\item[Raw input fingerprint (SHA256)] \path{2d219b5cb5947231612d4612e6f930e6d6ac740ea00d9cf2eaf9606c0bc772a2}
\item[Partition report SHA256] \path{9d0206908476cc6948b97307cad8ca3fae94e1d68aee0a3b45a69f697ade37f6}
\item[Boundary annotations SHA256] \path{85d3113e4f9e6fd7f6111b97f3d3b8ee14fc21392cfeeb1084181ad8c74306b2}
\item[Expanded review packets SHA256] \path{c76f8dc18e430e5dfb4215ec33749689608afd12213e33cd48392f14c8fa6df5}
\item[Expanded boundary annotations SHA256] \path{370d2fd4665ea0ce4b2e9df7905903bba0969ef310d0162c76dcdb348d85dc69}
\item[Expanded subject/admission review SHA256] \path{85845858aee76cc33bec0da699614913f86ad0aa84650dbf6c4d8566f7982191}
\end{description}
The annotations are exploratory single-assistant judgments, not independent
human labels. Future manuscript updates must preserve original run identities
and explain changes in input, annotation, parameters or implementation.
\end{document}
